%--------------------------------------------------------------------------------------
%------------------------------------- preambule --------------------------------------
%--------------------------------------------------------------------------------------
\documentclass[12pt]{ULojpv}

% Chargement des packages supplementaires
\usepackage[utf8]{inputenc}
\usepackage[T1]{fontenc}
\usepackage[french]{babel}
\usepackage{pifont}
\usepackage{comment}


% Definitions des parametres de l'en-tete
\Cours{GEL--1001 Design I (méthodologie)}             % Nom du cours
\NumeroEquipe{2}                                     % Numero de l'equipe
\NomEquipe{NordFaune}                               % Nom de l'equipe
\Objet{Procès-verbal}                                 % Nom du document
\SujetRencontre{Organisation et planification}             % Sujet de la rencontre
\DateRencontre{2026/01/29}                            % Date de la rencontre
\LocalRencontre{En ligne (Teams)}                            % Local de la rencontre
\HeureRencontre{8h30--9h20}                          % Heure de la rencontre


%--------------------------------------------------------------------------------------
%--------------------------------- corps du document ----------------------------------
%--------------------------------------------------------------------------------------
\begin{document}
\entete
\begin{enumerate}

% nouveau point
\item \textbf{Ouverture de la réunion}
Heure: 8h27

% nouveau point
\item \textbf{Nomination ou confirmation du président et du secrétaire}

\begin{tabular}{@{}ll}
   Président: Antoine Robert
   & Secrétaire: Louis-Charles Blais
\end{tabular}

% nouveau point
\item \textbf{Adoption de l'ordre du jour}

L'ordre du jour proposé est adopté à l'unanimité.

% nouveau point
\item \textbf{Lecture et adoption du procès-verbal de la réunion du 22 janvier 2026}

Il n'y avait pas de procès-verbal.

% nouveau point
\item \textbf{Affaires découlant du procès-verbal}

Il n'y avait pas de procès-verbal.

\begin{comment}
\begin{enumerate}

\item Affaire \#1

\begin{enumerate}

\item Sous-affaire \#1

Text here

\item Sous-affaire \#2

Text here

\end{enumerate}

\item Affaire \#2

Text here

\end{enumerate}
\end{comment}

% nouveau point
\item \textbf{Points à traiter}

\begin{enumerate}

\item Retour sur les séminaires du 27 et 28 janvier

Comme suggéré lors des séminaires, les membres de l'équipe ont partagé leur compréhension de l'objectif du client. Cela a permis de clarifier les demandes du client et d'orienter l’équipe dans la bonne direction.


\item Révision du Rapport de projet V0

Lorsque le projet a été comparé aux exemples fournis, on a trouvé différents problèmes : le nom du projet est incomplet, les noms ne sont pas en ordre alphabétique, il y a quelques fautes de français. Il a été décidé que chacun corrigerait ces erreurs de son côté immédiatement après la rencontre.


\item Organisation des futures remises

\begin{enumerate}

\item Documents de gestion G04

Puisque les membres de l'équipe ont eu de la difficulté avec la connexion au compte overleaf, l'équipe a décidé de transférer le travail vers GitHub. Toute l’équipe était d’accord sur l'importance de s'assurer que tous les membres se familiarisent avec les logiciels requis, soit Git, Youtrack et LaTeX. 


\item Rapport de projet - version 1

Les tâches de rédaction du projet n'ont pas encore été distribuées, afin de prioriser l'apprentissage des logiciels. Ce point sera discuté prochainement.

\end{enumerate}

\end{enumerate}


% nouveau point
\item \textbf{Divers}

Les personnes déjà familiarisées avec GitHub expliquent son utilisation aux autres membres. Antoine donne un tutoriel pour l’installation de Git sur Mac.


% nouveau point
\item \textbf{Répartition des tâches}

\begin{enumerate}

\item Ordre du jour, procès-verbal et diagramme de Gantt: L'ensemble du G04 devrait être transféré sur Github d'ici mercredi matin, afin de pouvoir réviser le contenu. Kayanne s'occupera de l'ordre du jour, Louis-Charles du procès-verbal et Mathis du diagramme de Gantt. La personne responsable de la remise du projet sera attribuée dans le groupe teams.

\end{enumerate}


% nouveau point
\item \textbf{Évaluation de la réunion}

La réunion a été efficace et productive. Toutes les questions et tous les doutes des membres ont été abordés durant la réunion. Globalement, l'équipe attribue une note de 9/10.


% nouveau point
\item \textbf{Date, heure, lieu et objectif de la prochaine réunion}

\begin{tabular}{@{}lll}
   Date: 2026/02/05
   & Heure: 08h30
   &  Lieu: En ligne (Teams)
\end{tabular}
\par
Discussion du rapport de projet V1


% nouveau point
\item \textbf{Fermeture de la réunion}

Heure: 09h21


% nouveau point
\item \textbf{Étaient présents}

\begin{dinglist}{"33}
   \item Antoine Robert
   \item Chelcie Laura Emmanuela Sonon (Partiellement)
   \item Félix-Antoine Guimont
   \item Kayanne Loudiyi
   \item Louis-Charles Blais
   \item Louise Geneviève Lesser
   \item Mathis Carrier
   \item Phanuel Kouadio
\end{dinglist}

\end{enumerate}

\end{document}

